% RESEARCH_DRAFT_SHA256: 774576d16febd46356d75e59315c7cff695657acd6b04e87d76a2129262894f2
% Layout and source-reference wrappers expanded for the web renderer.
\documentclass[11pt]{article}
\usepackage[margin=25mm]{geometry}
\usepackage{amsthm}
\usepackage{fontspec,unicode-math}
\setmainfont{Latin Modern Roman}
\setmathfont{Latin Modern Math}
\usepackage{amsmath,amssymb,mathtools}
\usepackage{xurl,hyperref}
\hypersetup{colorlinks=true,urlcolor=blue,linkcolor=blue}
\setcounter{secnumdepth}{0}
\setlength{\parindent}{0pt}
\setlength{\parskip}{5pt}
\setlength{\emergencystretch}{3em}
\newcommand{\sref}[2]{\hyperlink{source:#1:#2}{[S#2]}}
\newcommand{\cataloguescope}[1]{\paragraph{Source question.}#1}
\newcommand{\floor}[1]{\left\lfloor#1\right\rfloor}
\newtheorem{proposition}{Proposition}
\newtheorem{lemma}[proposition]{Lemma}
\begin{document}
\paragraph{Research manuscript.} OpenAI \texttt{gpt-6-astra}, reasoning \texttt{ultra}; 11 October 2026.
\paragraph{Preserved source body.} The complete catalogue mathematical text and sources follow with layout and source links expanded, including source qualifications and its historical preparation notes. New research begins at the section headed Research attempt.

\section{Moment-SOS Convergence for PDE-Constrained Variational Problems}\label{statement}

\subsection{Definitions and mathematical statement}

\paragraph{Shared notation.}Standard number systems are denoted by $\mathbb N,\mathbb Z,\mathbb Q,\mathbb R,\mathbb C$; index ranges, computational models and specialized notation follow the cited problem source and the definitions below.


A PDE-constrained variational problem optimizes a functional over functions satisfying a partial differential equation. Finite elements approximate those functions in finite-dimensional spaces; moment-SOS relaxations then replace finite polynomial constraints by increasingly strong positivity conditions.

A polynomial is a sum of squares (SOS) if it is $\sum_j q_j^2$ for polynomials $q_j$. Moment-SOS hierarchies impose positivity through increasingly large truncated moment matrices and corresponding degree-bounded SOS certificates.
\paragraph{Statement recorded in the supplied source.}
Extend the convergence analysis of finite-element moment-sum-of-squares relaxations beyond variational problems with a unique global minimizer and boundary constraints, to variational problems subject to partial differential equation constraints. \href{https://doi.org/10.4171/owr/2024/35}{[S2]}
\paragraph{Scope qualification.}
The PDE class, boundary conditions, admissible function space, and relaxation construction are source-specific. The target explicitly extends beyond unique minimizers.
\paragraph{Source status.}
The supplied catalogue labels this record open. Its source-status metadata was checked against the supplied snapshot on 2026-10-11. This import does not independently establish current open status or mathematical correctness. The catalogue literature assessment is dated 2026-08-22.
\subsection{Short English statement}
Prove convergence of the cited finite-element moment method for the broader PDE-constrained setting.
\subsection{Sources}
\begin{itemize}
\item[S1] Ulam AI and the UnsolvedMath Contributors, UnsolvedMath: A Curated Collection of Open Mathematics Problems, catalogue entry. Catalogue curation and metadata: CC BY 4.0. The supplied source question is retained; definitions and the English summary are editorial additions.
\url{https://huggingface.co/datasets/ulamai/UnsolvedMath}.
\item[S2] Polynomial Optimization for Nonlinear Dynamics: Theory, Algorithms and Applications (2025), Oberwolfach Reports, DOI 10.4171/owr/2024/35, pp. 1975--2032
\url{https://doi.org/10.4171/owr/2024/35}.
\textit{Underlying Oberwolfach Reports material retains its source-specific licence, including CC BY-SA 4.0 where applicable; see THIRD\_PARTY\_NOTICES.md.}
\end{itemize}
\paragraph{Preparation and review.}
Original question attributed to the cited source and UnsolvedMath contributors. Codex (GPT-6; exact serving revision not exposed) assisted with definitions, English summaries, source inspection and formatting. No new human mathematical authorship or independent proof verification is asserted. Catalogue identity checked against the complete supplied registry and existing repository statements; known flagged repetitions and obvious identical conjectures excluded. Full mathematical-equivalence/literature review has not been performed. Statement preparation and final review followed contribution prompts C and D; automated checks verify structure and identity, not open status.
\label{end:statement}

\section{Research attempt}
\subsection{Literature, source context and scope}
The full original catalogue body and source qualifications are preserved above. The relevant primary workshop passages were read in the cached publisher report. Historical claims and subsequent attributed results are distinguished from the mathematical work below. A restricted example, standard lemma or known theorem is not counted as new progress on the entire original target. Missing assumptions and repeated or rhetorical extracts are not silently turned into different solved problems.

The primary passage distinguishes discretization followed by moment relaxation from certified lower-bound approaches. It reports convergence to a minimizer under uniqueness and only boundary constraints, then asks about nonuniqueness and further PDE constraints. This attempt tests whether optimal relaxed moments can simply be averaged to remove uniqueness.
\subsection{A PDE-constrained example with exact moments and failed mean extraction}
On $(0,1)$ take $u\in H^1(0,1)$, weak constraint $u'=0$ and $|u|\le1$ almost everywhere, with no imposed boundary values, and minimize
\[
 F(u)=\int_0^1(u(x)^2-1)^2\,dx.
\]
The derivative constraint implies $u=c$ almost everywhere: an $H^1$ function has an absolutely continuous representative with $u(x)-u(0)=\int_0^xu'$. Thus the problem is exactly $\min_{-1\le c\le1}(c^2-1)^2$, with two minimizers $c=\pm1$ and optimum zero. Constant finite elements already represent both solutions, so no mesh error is involved.
\begin{proposition}
The order-two moment relaxation has an optimal feasible moment vector whose first-moment field is feasible but has objective one, strictly above the true optimum.
\end{proposition}
\begin{proof}
Set $(y_0,y_1,y_2,y_3,y_4)=(1,0,1,0,1)$, the moments of $(\delta_{-1}+\delta_1)/2$. Its moment and localizing matrices are
\[
 M_2(y)=\begin{pmatrix}1&0&1\\0&1&0\\1&0&1\end{pmatrix}\succeq0,
 \qquad M_1((1-c^2)y)=\begin{pmatrix}0&0\\0&0\end{pmatrix}\succeq0.
\]
Indeed the first quadratic form is $(v_0+v_2)^2+v_1^2$. The relaxation minimizes $y_4-2y_2+y_0$, which is nonnegative for every feasible moment matrix because it is the moment functional applied to the square $(c^2-1)^2$. Our vector has value zero and is optimal. Its first moment is zero, so the mean reconstruction is $u(x)=0$, satisfying the PDE and bound but yielding $F(u)=1$.
\end{proof}
The same measure supplies compatible optimal moments at every higher order, so raising relaxation order alone does not cure this particular extraction rule. On the other hand, selecting either support atom gives a true minimizer. The obstruction is specifically to replacing a measure supported on minimizers by its barycenter for a nonconvex objective.
\section{Verification and remaining gaps}
Exact rational checks verify the two matrices, the zero relaxed objective and the reconstructed objective one. This does not refute convergence in value, convergence to measures supported on minimizers, or a valid atom-selection theorem. General PDE discretization, compactness, relaxation exactness and a justified extraction rule without uniqueness are all still missing.

\section{References and source attribution}
Primary workshop report: \url{https://doi.org/10.4171/owr/2024/35}. Relevant passages read: OWR 2024/35, printed pages 1991--1992, Fantuzzi on polynomial optimization for variational problems.

All retrievals were checked on 11 October 2026. Original question attribution remains in the preserved Sources section. External theorems are attributed and not claimed as new discoveries or fully verified merely by citation.


\paragraph{Integration review.} Codex (GPT-6; exact serving revision not exposed) reviewed the written arguments and expanded layout and source-reference wrappers. Repository metadata and the production web renderer were checked using Python 3.13 and JavaScript with MathJax 3.2.2. This model-assisted review is not independent human proof verification; no PDF compilation or formal proof checking is asserted.

\section{Completion Estimate}
\noindent\textbf{COMPLETION: 0\%}
No substantiated new progress toward the entire original target is claimed. The worked results and failed route are delimited above; a known result, illustrative model, or auxiliary obstruction is not counted as original-target completion.

\section{Final status and attribution}
\textbf{Attempt status: unresolved.} This describes this research attempt, not a new ruling on the source catalogue's status or on external literature claims.
The mathematics and this self-audit were generated by OpenAI \texttt{gpt-6-astra}, reasoning \texttt{ultra}, on 11 October 2026. Codex (GPT-6; exact serving revision not exposed) orchestrated the batch, read and reviewed the written mathematical arguments, independently reran the finite checks, and contributed formatting/repository review. No human mathematical authorship, independent human proof checking, or formal verification is claimed. The historical preparation paragraphs in the reproduced catalogue body concern that body only.

Contribution prompt A was applied to the complete definition and supplied source record, followed by the adversarial checks of prompt B. Local reproduction files are \texttt{batch026/checks.py}, \texttt{computations.json}, and the per-ID audit JSON. They were not public links at preparation time. Numerical checks are finite and do not prove unrestricted results. No LaTeX engine was found locally; compilation was not performed by this research agent.

\end{document}
